% =====================================================================
%  Extended captions for the six COE validation panels.
%  Included via \input from the main paper inside a figure float.
% =====================================================================

\newcommand{\captionCOEpanelone}{%
\textbf{Time-Count Identity and count linearity.}
(a) Recovered reference frequency $\hat f = M(t)/t$ across $50$ trials at
$f_\mathrm{ref} = 10^6\,\mathrm{Hz}$; the dashed red line marks the reference
frequency. The measurement and the operational definition $t = \Mcount/f$
agree to within sub-permille precision across the full range.
(b) Per-trial relative error $|\hat f - f_\mathrm{ref}|/f_\mathrm{ref}$ on a
log scale; all trials lie at least three orders of magnitude below the
$10^{-3}$ tolerance line.
(c) Linearity test: $M(\op_1 \cup \op_2)$ versus $M(\op_1) + M(\op_2)$
across $50$ randomly chosen non-overlapping operation pairs; points fall
exactly on the identity line, confirming that the count is additive
on disjoint operation unions.
(d) Three-dimensional surface of the conversion law $t = Q/f$ over
$Q \in [1, 200]$ and $\log_{10} f \in [3, 7]$. The surface visualises
how a single underlying $Q$ projects to different $(t, M)$ pairs through
the kernel's reference frequency.}

\newcommand{\captionCOEpaneltwo}{%
\textbf{Three independent routes to operation weight.}
(a) Route I (residue): $Q_\mathrm{I}$ measured by an instrumented
decision-counter against the ground-truth weight $Q_\mathrm{true}$; the
match is exact across $30$ trials with $Q_\mathrm{true} \in [1, 10^4]$.
(b) Route II (confinement): $Q_\mathrm{II}$ measured by an
address-space-cell-counter against $Q_\mathrm{true}$; same exact match.
(c) Route III (negation fixed point): $Q_\mathrm{III}$ measured by an
iterated-negation-depth instrument against $Q_\mathrm{true}$; same exact match.
(d) Three-dimensional point cloud of the triples
$(Q_\mathrm{I}, Q_\mathrm{II}, Q_\mathrm{III})$ for $30$ trials; every point
lies on the diagonal $Q_\mathrm{I} = Q_\mathrm{II} = Q_\mathrm{III}$
(red dashed line), demonstrating that the three computationally distinct
instruments converge to a single underlying invariant.}

\newcommand{\captionCOEpanelthree}{%
\textbf{Three-route equivalence and Mass-Time-Identity-Count.}
(a) Pairwise differences $Q_\mathrm{I} - Q_\mathrm{II}$,
$Q_\mathrm{II} - Q_\mathrm{III}$, $Q_\mathrm{I} - Q_\mathrm{III}$ across $50$
trials; all curves lie identically on zero, confirming Theorem~\ref{thm:three-route}
to machine precision.
(b) $Q_\mathrm{I}$ versus $Q_\mathrm{II}$ coloured by $Q_\mathrm{III}$
intensity; the colour gradient runs along the diagonal,
demonstrating that all three quantities co-vary perfectly.
(c) Mass-Time-Identity-Count Equivalence: $t = Q/f$ recovered both directly
(blue) and from $Q_\mathrm{from\,t}/f$ (amber crosses), showing that the
four-way unit conversion is consistent in both directions.
(d) Three-dimensional locus of $(Q, M, t)$ points, lying on the
one-parameter curve $Q = M$, $t = Q/f$. The MTIC theorem states that this
curve is the only locus a kernel can occupy: any divergence between
$Q$, $M$, and $tf$ is a measurement artefact, not a degree of freedom.}

\newcommand{\captionCOEpanelfour}{%
\textbf{Reproducibility, sliding-endpoint, and monotone log growth.}
(a) Reproducibility: $1000$ runs of an identical operation produce a single
distinct output, weight, and hash; all three identical-flags read $1$,
confirming the reproducibility property.
(b) Sliding-endpoint theorem: the baseline configuration is reproducible
(green bar at $1$), and externally truncating the kernel's accumulated count
breaks reproducibility (amber bar at $1$), confirming
Theorem~\ref{thm:sliding}: reproducibility $\iff$ irreversibility of $\Mcount$.
(c) Rewind-as-Forward Principle: $\Mcount$ after rollback (y-axis) versus
$\Mcount$ before rollback (x-axis) across $30$ trials; every point lies above
the equality line, demonstrating that rollback strictly increases the count.
(d) Monotone log growth: 3D bars showing $\Delta \Mcount$ at each operation
step over the first $50$ ops of a $1000$-op workload mixing forwards and
rollbacks. Every bar has nonnegative height, demonstrating that the log
grows monotonically by structural force, not by engineering choice.}

\newcommand{\captionCOEpanelfive}{%
\textbf{API substitutability and cell stability.}
(a) Two cascade chains $A$ and $B$, internally distinct, produce identical
outputs across $50$ trials; points lie exactly on the identity line.
The fixed-point semantics of the API boundary make the chains
operationally indistinguishable from any external observer.
(b) Cell-stability test: existing operations remain unaffected when a
cell-disjoint operation is added; all $30$ trials read $1$ (stable).
(c) Cell-collision detection: confusion matrix with $50$ trials
(half disjoint, half colliding additions). All decisions land on
the diagonal --- no false positives, no false negatives ---
demonstrating perfect collision detection.
(d) Disjoint cell ownership: 3D bar chart of operation weights $Q$
across cell indices and kernel identities. Each kernel owns a disjoint
subset of cells; the absence of vertical bars at non-owned cells
visualises the cell-stability invariant.}

\newcommand{\captionCOEpanelsix}{%
\textbf{Cross-architecture invariance.}
(a) The relation $t = Q/f$ holds simultaneously for four kernel
architectures (microkernel, monolithic, exokernel, unikernel),
each operating at its own reference frequency. Each architecture's points
follow the conversion law with no architecture-specific corrections.
(b) Reference-frequency hierarchy: bar chart on a logarithmic axis
showing the four architecture-local frequencies, spanning over a decade.
(c) Per-architecture recovery error $|Q - Q_\mathrm{recovered}|$ across
$10$ trials per architecture; all curves lie at the floor of the log
plot ($+10^{-3}$ for visualisation), confirming that the four-way
unit conversion is exact independent of $f$.
(d) Three-dimensional view of $(Q, t)$ stratified by architecture; the
four point clouds lie on architecture-specific lines through the origin,
visualising that $f$ is the only architecture-dependent quantity in the
MTIC framework.}
